% ---
% Capitulo de Conclusao
% ---

\chapter{Conclusão\label{cha:conclusion}}

Este trabalho tornou a construção do \ac{CORE-SG} executável em múltiplos aceleradores, tratando seus dois grafos constituintes como problemas separados: uma decomposição em instâncias independentes para o grafo de vizinhança, e uma estimativa por pontes anti-hub, dimensionada independentemente de $n$, para a árvore geradora.\\

O problema atacado, enunciado na Seção \ref{sec:problem-statement}, era de viabilidade. A construção original mantém as duas metades em um único espaço de memória, de modo que o maior conjunto processável fica determinado pela capacidade de um dispositivo, e acrescentar aceleradores não altera esse limite. As quatro questões de pesquisa da Seção \ref{sec:research-questions} decompõem esse problema, e o Capítulo \ref{cha:results} as responde com medições; a Seção \ref{sec:disc-summary} reúne as respostas.\\

Antes da decomposição, o maior conjunto processável era determinado pela memória de um único dispositivo. Depois dela, a residência por dispositivo é determinada por um parâmetro que o usuário escolhe, e a quantidade de aceleradores determina o tempo de execução. O limite de escala não desaparece: ele passa do dispositivo para a memória do hospedeiro, que continua a crescer com $n$. A Seção \ref{sec:res-scale} exibe a consequência: construções de $10^7$ a $5 \times 10^7$ pontos mantêm a residência por dispositivo em torno de 293\,MB, enquanto a construção não decomposta sobre os mesmos dados excede a capacidade do acelerador a partir de $10^7$.\\

A metade de árvore geradora admite substituição por uma estimativa cujo custo não é governado pela quantidade de pontos. As pontes anti-hub partem de uma estatística que o próprio grafo de vizinhança já contém, o grau reverso normalizado, e o conjunto de vértices que elas abrangem é dimensionado pela dimensionalidade e pela quantidade de arestas. As árvores assim obtidas reproduzem o peso total da árvore exata com erro máximo de 0,024\% nos conjuntos de $n = 4.000$ em que a referência exata é calculável, quase todos sintéticos e construídos para fragmentar o grafo de vizinhança, e a Seção \ref{sec:res-connectivity} quantifica, em partes por milhão, o quanto a metade de árvore geradora contribui nos conjuntos reais, em que o grafo de vizinhança já é conexo.\\

O custo da decomposição está localizado e foi medido. O particionamento deixa pares não comparados, e recuperá-los exige um passo de reparo que executa no hospedeiro e responde por cerca de quatro quintos de uma construção em escala. Quatro estratégias de redução foram implementadas e medidas, e nenhuma delas encurtou o passo a qualidade constante; o que desloca parte de seu trabalho para os aceleradores é a fração de sobreposição, que a Seção \ref{sec:res-overlap} caracteriza. Essa é a limitação dominante do método, e a Seção \ref{sec:disc-limits} enumera as demais, junto com as condições em que as medições foram tomadas.\\

As direções de continuação estão na Seção \ref{sec:disc-future}. A mais imediata é a transposição da fase de reparo para o acelerador, que é onde o limite de escalabilidade atual reside. A Seção \ref{sec:fut-streams} registra os resultados preliminares sobre fluxos de dados, obtidos antes da consolidação do escopo atual, porque a abstração de instância se aplica à extensão de um grafo existente sem alteração estrutural, o que faz dessa a continuação natural da linha.\\

Três constatações deste trabalho não dependem do \ac{CORE-SG}. A primeira é que, em um sistema que combina hospedeiro e aceleradores, a grandeza que governa o projeto é a distribuição do custo entre os dois lados. A segunda é que a estatística de \textit{hubness}, tratada na literatura como um problema a mitigar, é um estimador de borda utilizável quando invertida e normalizada localmente. A terceira é que um grafo de vizinhança em alta dimensionalidade não admite partição cujos blocos tenham fronteiras pequenas, de modo que aceitar a perda e recuperá-la depois é uma resposta mais praticável que procurar uma partição melhor. A Seção \ref{sec:disc-insights} desenvolve as três.\\
